\section*{Opening: The Bridge}

This paper dives into the intersection of literacy and social justice by showing how communication systems can empower or exclude communities. In the BIO/STEM communication-access case, course literacy did not only mean reading and writing. It meant being able to decode requirements, accommodations, grading language, single-mode directions, team expectations, scientific communication norms, and institutional pathways. When a teacher uses their authority, knowledge, or influence to communicate through one-size-fits-all expectations, they can intentionally or unintentionally exclude neurodivergent, chronically ill, disabled, culturally and linguistically diverse, trauma-impacted, and otherwise vulnerable students.

The purpose of this paper is not to reduce the issue to one person or one course. The purpose is to use a documented adult learning experience as a case for analyzing a larger social justice problem: communication access injustice in education. The central question is this: if an adult graduate student with technical knowledge, documentation habits, advisor support, and self-advocacy skills can still be harmed by inaccessible communication, what happens to children and younger students who cannot defend themselves yet?

\placeholderfigure{communication-access-map}{Insert a conceptual figure showing the relationship among classroom communication, disciplinary literacy, accommodations, institutional navigation, and student access.}{This figure can later become a polished visual model of the paper's argument.}

\section*{Research Findings: Communication Access as a Social Justice Issue}

\subsection*{Identify the Topic}

The social justice issue I am focusing on is \textbf{communication access injustice for neurodivergent and chronically ill students in STEM education}. This issue sits directly at the intersection of literacy and social justice because literacy is not only the ability to read and write words. In school, literacy also means being able to access, decode, interpret, and respond to the communication systems that organize learning.

In a STEM classroom, those communication systems include syllabi, assignment instructions, grading policies, verbal updates, accommodation procedures, email exchanges, peer evaluation expectations, team communication platforms, scientific terminology, code documentation, file formats, and institutional support pathways. A student is not only learning biology, bioinformatics, computer science, or data analysis. The student is also learning how to read the course itself.

When those systems are clear, stable, accessible, and multimodal, literacy can empower students. It can help them understand expectations, ask better questions, use accommodations, collaborate with peers, and show what they know. But when those systems are vague, single-mode, inconsistent, hidden, or built around one assumed kind of learner, literacy becomes a gatekeeping system. Students who do not process, remember, communicate, or participate in that default way can be misread as careless, difficult, unprepared, irresponsible, or less capable.

That is why this issue is social justice, not just classroom preference. A teacher's communication choices decide who gets access to learning and who has to fight to be understood. Even when a teacher does not intend harm, their authority, disciplinary knowledge, and classroom structure can exclude students if communication is designed for only one kind of body, brain, language background, pace, memory system, or emotional response.

The BIO/STEM evidence gathered for this paper documents this issue through an adult experience in a bioinformatics course. However, the adult case is not the endpoint of the paper. It is the evidence site. The larger issue is how inaccessible communication affects neurodivergent, chronically ill, disabled, culturally and linguistically diverse, trauma-impacted, anxious, depressed, and otherwise vulnerable students across educational settings.

\subsection*{Why This Issue Matters}

This issue matters because communication determines access. If a student cannot reliably find, understand, remember, or clarify the language of a course, then the student cannot fully access the learning the course claims to offer. Inaccessible communication can affect far more than grades. It can affect physical health, executive functioning, mental health, peer relationships, academic confidence, trust in teachers, willingness to seek help, identity as a capable learner, and future participation in STEM.

The issue is historically relevant because schools have often organized learning around a narrow image of the successful student: someone who can sit still, remember oral instructions, use dominant academic language, respond quickly, communicate calmly under stress, tolerate ambiguity, and perform knowledge in approved formats. Students who do not match that image are often misread. Neurodivergent students may be seen as careless, disruptive, lazy, too intense, or disorganized. Chronically ill students may be seen as unreliable. Culturally and linguistically diverse and multilingual students may be seen as less prepared when the real issue is institutional language. Students with trauma, depression, or anxiety may be seen as avoidant or difficult when the communication environment itself feels threatening.

This becomes especially urgent in STEM because STEM courses can confuse rigor with rigidity. A course can be intellectually demanding without being inaccessible. But when instructors treat ambiguity, memory for oral directions, single-mode output, public correction, or emotional toughness as evidence of rigor, students with access needs are quietly sorted out. The problem is not that STEM should become easier. The problem is that access should be part of rigor.

This issue also matters in my field of study because bioinformatics and computer science are deeply literate fields. They require students to read scientific instructions, interpret computational tasks, understand data files, communicate through code, document processes, coordinate digital teamwork, translate technical thinking into reports, and navigate institutional policies. If a course treats these literacies as obvious instead of teachable and accessible, it can exclude students even when they have the intellectual capacity to do the work.

The local community connection is also important. In Literacy Learning as Social Practice and in my teaching placement, I have seen how much students, especially children, rely on clear communication, modeling, visuals, repetition, embodied practice, peer explanation, and emotionally safe clarification. Children often cannot document harm, email advisors, cite accommodations, or explain that a classroom expectation is inaccessible by design. They may simply shut down, act out, become overwhelmed, or internalize the message that they are the problem.

That is the reason this issue needs attention. If an adult graduate student with documentation habits, technical expertise, advisor support, and self-advocacy skills can still be harmed by inaccessible communication, then children and less-resourced students are at even greater risk. They should not have to become experts in institutional navigation before they are allowed to be learners.

\subsection*{Relevant Sources and Evidence Scope}

The research for this paper draws from four bodies of evidence: a curated BIO/STEM communication-access case file, literacy and social justice course readings, teaching-placement observations, and a present-day accessibility initiative. Together, these sources connect a documented adult learning experience to broader questions of literacy, power, disability, disciplinary participation, and educational design.

The \textbf{BIO/STEM communication-access case file} serves as the central evidence source for the adult case. It brings together course communication, selected classroom and team records, advisor-supported institutional navigation, and examples of scientific and technical communication. The case focuses on how access barriers emerged through course requirements, grading procedures, accommodation communication, team coordination, and disciplinary expectations. Examples include undocumented or unstable procedural expectations, grading language that required interpretation beyond the written policy, health-related communication that increased the burden on a chronically ill student, accommodation procedures experienced through compliance limits rather than inclusive design, collaborative project expectations without sufficient communication norms, and specialized scientific terminology treated as self-evident rather than explicitly supported.

The \textbf{literacy and social justice readings} provide the conceptual framework for interpreting the BIO/STEM case. These readings position literacy as more than a technical skill. They connect literacy to identity, disciplinary participation, criticality, language power, access, and joy. This framework matters because the paper is not only asking what happened in one course. It is asking how classroom communication can position students as capable or incapable, included or excluded, legitimate or difficult.

The \textbf{teaching-placement observations} provide the K--12 bridge. They connect the adult BIO/STEM case to children who often demonstrate understanding through action, tools, collaboration, movement, design, peer explanation, and repeated guided practice. These observations matter because children usually cannot document harm, email advisors, cite accommodations, or explain that an expectation is inaccessible by design. Teaching placement shows why communication access must be designed before harm occurs, not repaired only after students have already been misread.

The \textbf{present-day accessibility initiative}, Universal Design for Learning, especially CAST UDL Guidelines 3.0, supports the move from documentation to redesign. This initiative belongs mainly in the later section connecting the issue to the present, but it is included here as part of the research direction because it offers a current framework for designing learning environments that anticipate learner variability rather than treating difference as an exception \citep{cast2024udl}.

Together, these sources establish the research foundation for the paper. They identify the topic, show why it matters, and prepare the critical analysis that follows: literacy is not only whether students can read words on a page. Literacy is whether students can access the communication systems that decide grades, participation, accommodations, belonging, and recognition.

\section*{Critical Analysis: Literacy as Communication Power}

Literacy shapes how communication access injustice is understood because the issue is often hidden behind language that sounds neutral. Words like \textit{rigor}, \textit{professionalism}, \textit{participation}, \textit{responsibility}, \textit{reasonable accommodation}, and \textit{student expectations} can appear fair on the surface while still placing uneven burdens on students whose bodies, brains, languages, or trauma histories do not match the assumed classroom norm. In this way, literacy does not simply communicate a social justice issue; it can also disguise the issue.

In the BIO/STEM communication-access case, the problem was not only that individual instructions were unclear. The deeper issue was that course communication operated through a set of literacy demands that were treated as obvious: knowing which expectations were official, knowing how to interpret grading language, knowing when a classroom direction carried the force of policy, knowing how to translate illness into institutional procedure, knowing how to use accommodations without appearing difficult, and knowing how to maintain team participation while managing chronic illness. These are not minor communication details. They are access conditions.

\subsection*{Language, Tone, and Authority}

The language of a course carries institutional power because students rarely meet course communication as equals. A professor's wording can determine what counts as responsible, late, complete, professional, respectful, or valid. When expectations are communicated through unstable, single-mode, or implied channels, the student must do more than complete the assignment. The student must interpret the system.

This matters because the same sentence can have different consequences depending on who receives it. A student without health barriers may experience an unclear requirement as an inconvenience. A neurodivergent student may experience it as an executive-function crisis. A chronically ill student may experience it as another demand layered on top of symptoms, medication effects, and disclosure fatigue. A child may not experience it as communication ambiguity at all; they may simply experience shame, shutdown, or punishment.

The BIO/STEM case shows how communication can become exclusionary even when it is framed as ordinary course management. Procedural expectations communicated through a single mode can become inaccessible when students rely on stable written records. Grading language requiring interpretation beyond the written policy can turn assessment into policy decoding. Health-related communication can shift responsibility back to the student during moments when the student has the least capacity. Accommodation language can become a compliance checkpoint rather than an access design. Team communication can become another barrier when collaboration tools exist without explicit norms.

The issue, then, is not whether communication exists. The issue is whether communication is accessible, stable, interpretable, and humane.

\subsection*{Terminology as Access or Exclusion}

Specific terms can exclude students when they are used as if everyone shares the same institutional, disciplinary, or disability literacy. In this case, terms such as \textit{Flex Plan}, \textit{DSO}, \textit{Ombuds}, \textit{peer evaluation}, \textit{formative feedback}, \textit{participation}, \textit{professional communication}, and \textit{rigor} are not simple vocabulary words. They are gatekeeping terms when students are expected to understand their consequences without explicit explanation.

Scientific and technical terminology can create a similar barrier. In bioinformatics and computer science, students are expected to understand not only concepts, but also formats, workflows, documentation practices, computational tools, and disciplinary conventions. Specialized language is necessary in STEM, but it becomes exclusionary when it is treated as self-evident rather than taught, modeled, or connected to clear expectations.

This distinction matters. The goal is not to remove specialized language from STEM. That would weaken disciplinary learning. The goal is to make specialized language accessible. Students should be able to learn the language of the field without also having to decode hidden classroom rules, unclear grading procedures, or institutional processes during moments of crisis.

\subsection*{Audience Literacy and Unequal Burden}

Different audiences bring different kinds of literacy to this issue. A graduate student may have advanced scientific literacy but still struggle with institutional ambiguity, accommodation procedures, or unclear grading language. A neurodivergent student may have strong analytical literacy but need stable written expectations, predictable channels, and explicit sequencing. A chronically ill student may understand the work but need communication structures that account for variable capacity. A culturally and linguistically diverse or multilingual student may understand concepts but be excluded by unexplained institutional language or dominant academic tone. A student with trauma, anxiety, or depression may experience public correction, uncertainty, or conflict as threat rather than feedback.

Children are at even greater risk because they usually lack the institutional language adults can use to defend themselves. They may not know how to say that a classroom expectation is inaccessible, that a direction was too vague, that they needed a visual model, that public correction made them shut down, or that they understood the idea but could not perform it in the required format. Instead, their behavior may be misread. They may be labeled inattentive, oppositional, unmotivated, careless, or disruptive.

This is where the adult BIO/STEM case becomes a warning. If an adult with documentation practices, technical expertise, advisor support, and self-advocacy skills can still be harmed by inaccessible communication, then children and less-resourced students face even greater danger. They are less likely to have the language, power, or protection needed to challenge the system. For them, communication access cannot be optional. It must be designed before harm occurs.

\subsection*{Literacy as the Mediator of Social Justice}

The course readings on literacy, identity, disciplinary learning, and criticality help explain why this is a social justice issue. Literacy is not neutral. It positions students. It can recognize students as capable, or it can make their needs appear excessive. It can invite participation, or it can make participation dependent on decoding hidden rules. It can support identity, or it can force students to mask, over-document, or repeatedly justify their existence in the classroom.

Through this lens, communication access is not a soft addition to instruction. It is part of educational justice. A classroom that values only one way of processing, remembering, speaking, writing, participating, or demonstrating knowledge is not rigorous in a meaningful sense. It is selective. It rewards students who already fit the expected literacy profile and penalizes students whose literacies appear in different forms.

A more just classroom does not lower standards. It clarifies standards. It makes expectations visible. It teaches disciplinary language explicitly. It repeats important changes in writing. It treats accommodations as a starting point for design, not a ceiling. It recognizes that students demonstrate understanding through multiple literacies: written, oral, visual, technical, embodied, collaborative, emotional, and institutional.

This is the bridge to the recommendations that follow. The solution is not only better individual kindness. The solution is better communication design. Literacy can exclude, but it can also empower when educators design systems that allow students to understand expectations, access support, participate with dignity, and show what they know without having to break themselves first.
